\documentclass[10pt,a4paper]{amsart}
\usepackage{amsmath,amssymb,amsthm,graphicx}
% For commutative diagrams we use Xy-pic
\usepackage[all]{xy}

% We use 2cell for 2-commutative diagrams.
\xyoption{2cell}
% Theorem environments.
%
\theoremstyle{plain}
\newtheorem{theorem}[subsection]{Theorem}
\newtheorem{proposition}[subsection]{Proposition}
\newtheorem{lemma}[subsection]{Lemma}

\theoremstyle{definition}
\newtheorem{definition}[subsection]{Definition}
\newtheorem{example}[subsection]{Example}
\newtheorem{exercise}[subsection]{Exercise}
\newtheorem{situation}[subsection]{Situation}

\theoremstyle{remark}
\newtheorem{remark}[subsection]{Remark}
\newtheorem{remarks}[subsection]{Remarks}

\numberwithin{equation}{subsection}

% Macros
%
\def\lim{\mathop{\mathrm{lim}}\nolimits}
\def\colim{\mathop{\mathrm{colim}}\nolimits}
\def\Spec{\mathop{\mathrm{Spec}}}
\def\Hom{\mathop{\mathrm{Hom}}\nolimits}
\def\Ext{\mathop{\mathrm{Ext}}\nolimits}
\def\SheafHom{\mathop{\mathcal{H}\!\mathit{om}}\nolimits}
\def\SheafExt{\mathop{\mathcal{E}\!\mathit{xt}}\nolimits}
\def\Sch{\mathit{Sch}}
\def\Mor{\mathop{\mathrm{Mor}}\nolimits}
\def\Ob{\mathop{\mathrm{Ob}}\nolimits}
\def\Sh{\mathop{\mathit{Sh}}\nolimits}
\def\NL{\mathop{N\!L}\nolimits}
\def\CH{\mathop{\mathrm{CH}}\nolimits}
\def\proetale{{pro\text{-}\acute{e}tale}}
\def\etale{{\acute{e}tale}}
\def\QCoh{\mathit{QCoh}}
\def\Ker{\mathop{\mathrm{Ker}}}
\def\Im{\mathop{\mathrm{Im}}}
\def\Coker{\mathop{\mathrm{Coker}}}
\def\Coim{\mathop{\mathrm{Coim}}}

% Boxtimes
%
\DeclareMathSymbol{\boxtimes}{\mathbin}{AMSa}{"02}

%
% Macros for moduli stacks/spaces
%
\def\QCohstack{\mathcal{QC}\!\mathit{oh}}
\def\Cohstack{\mathcal{C}\!\mathit{oh}}
\def\Spacesstack{\mathcal{S}\!\mathit{paces}}
\def\Quotfunctor{\mathrm{Quot}}
\def\Hilbfunctor{\mathrm{Hilb}}
\def\Curvesstack{\mathcal{C}\!\mathit{urves}}
\def\Polarizedstack{\mathcal{P}\!\mathit{olarized}}
\def\Complexesstack{\mathcal{C}\!\mathit{omplexes}}
% \Pic is the operator that assigns to X its picard group, usage \Pic(X)
% \Picardstack_{X/B} denotes the Picard stack of X over B
% \Picardfunctor_{X/B} denotes the Picard functor of X over B
\def\Pic{\mathop{\mathrm{Pic}}\nolimits}
\def\Picardstack{\mathcal{P}\!\mathit{ic}}
\def\Picardfunctor{\mathrm{Pic}}
\def\Deformationcategory{\mathcal{D}\!\mathit{ef}}
\usepackage{cleveref}
\crefname{originallemma}{original lemma}{original lemmas}
\Crefname{originallemma}{Original lemma}{Original lemmas}
\makeatletter
\newlabel{lemma-geometric-cap}{{02SQ}{source}{}{sourcepointer-1}{}}
\newlabel{lemma-geometric-cap@cref}{{[originallemma][0][]02SQ}{[1][0][]source}{}{}{}}
\makeatother
\crefname{originallemma}{original lemma}{original lemmas}
\Crefname{originallemma}{Original lemma}{Original lemmas}
\makeatletter
\newlabel{situation-setup}{{02QL}{source}{}{sourcepointer-2}{}}
\newlabel{situation-setup@cref}{{[originallemma][0][]02QL}{[1][0][]source}{}{}{}}
\makeatother
\title{Pushing forward powers of the first Chern class of the tautological bundle on a projective bundle}
\author{The Stacks Project Authors}
\date{stacks-project:chow.tex:local-snapshot-20261012; GNU FDL 1.2 or later}
\begin{document}
\maketitle
\noindent\textit{Editorial note.} All delivered mathematical text is the author's own, copied from the pinned Stacks Project chapter file chow.tex; no statement or proof step was rewritten, reordered or removed.\par\medskip
\noindent\textit{Editorial note.} Preamble provenance: the pinned chapter file opens with \textbackslash{}input\{preamble\} and carries no preamble text of its own, so the project's single Stacks-wide preamble staged at sources/stacks-chow/preamble.tex is excerpted verbatim as two spans, lines 14-18 (Xy-pic setup) and lines 154-217 (the native theorem declarations and mathematical macros). Nothing was added, removed or redefined; the wrapper class and the amsmath/amssymb/amsthm/graphicx base packages are supplied by the assembler.\par\medskip
\noindent\textit{Editorial note.} Same-chapter pointers are located at their \textbackslash{}label lines inside the pinned chapter file; cross-chapter pointers are located in the project's own tag list staged in the snapshot (sources/stacks-chow/tags.txt), at the verbatim lines carrying their canonical tags.\par\medskip
\noindent\textit{Editorial note.} The author's closing sentence at source line 6927, 'We omit the verification that these arguments imply the result for a general cycle alpha = sum n\_j [Z\_j].', is retained verbatim at the end of the counted proof. The construction of the rational sections, the iteration of the geometric case and the pushforward computation are all inside the span; only the passage from an integral class to a locally finite sum of them is left to the reader.\par\medskip
\noindent\textit{Allowed standard prerequisites (editorial).}\begin{itemize}
\item The geometric case of capping with a first Chern class. Same chapter chow.tex, label lemma-geometric-cap, invoked at proof line 6924 and used repeatedly from proof line 6916 to 6923. Identifies the successive caps c\_1(O\_P(1))\textasciicircum{}i cap pi\textasciicircum{}*[Z] with the cycle classes of the common zero schemes Z(s\_1) intersect ... intersect Z(s\_i) of the constructed rational sections.
\end{itemize}\medskip
\ifdefined\hypertarget\hypertarget{sourcepointer-1}{}\fi
\noindent\textit{Source reference.} Original lemma 02SQ: sources/stacks-chow/pinned-original-source.tex, source lines 4432--4432. Navigation pointer for \textbackslash{}ref\{lemma-geometric-cap\} used inside the counted span (line 6924); the argument that consumes it is retained in the span.\par\medskip
\ifdefined\hypertarget\hypertarget{sourcepointer-2}{}\fi
\noindent\textit{Source reference.} Original lemma 02QL: sources/stacks-chow/pinned-original-source.tex, source lines 1372--1372. Navigation pointer for \textbackslash{}ref\{situation-setup\} used inside the counted span (line 6869); the argument that consumes it is retained in the span.\par\medskip
\begin{lemma}
\label{lemma-cap-projective-bundle}
Let $(S, \delta)$ be as in Situation \ref{situation-setup}.
Let $X$ be locally of finite type over $S$.
Let $\mathcal{E}$ be a finite locally free $\mathcal{O}_X$-module
$\mathcal{E}$ of rank $r$. Let $(\pi : P \to X, \mathcal{O}_P(1))$
be the projective bundle associated to $\mathcal{E}$.
For any $\alpha \in \CH_k(X)$ the element
$$
\pi_*\left(
c_1(\mathcal{O}_P(1))^s \cap \pi^*\alpha
\right)
\in
\CH_{k + r - 1 - s}(X)
$$
is $0$ if $s < r - 1$ and is equal to $\alpha$ when $s = r - 1$.
\end{lemma}

\begin{proof}
Let $Z \subset X$ be an integral closed subscheme of $\delta$-dimension $k$.
Note that $\pi^*[Z] = [\pi^{-1}(Z)]$ as $\pi^{-1}(Z)$ is integral of
$\delta$-dimension $r - 1$.
If $s < r - 1$, then by construction
$c_1(\mathcal{O}_P(1))^s \cap \pi^*[Z]$
is represented by a $(k + r - 1 - s)$-cycle supported on
$\pi^{-1}(Z)$. Hence the pushforward of this cycle
is zero for dimension reasons.

\medskip\noindent
Let $s = r - 1$. By the argument given above we see that
$\pi_*(c_1(\mathcal{O}_P(1))^s \cap \pi^*\alpha) = n [Z]$
for some $n \in \mathbf{Z}$. We want to show that $n = 1$.
For the same dimension reasons
as above it suffices to prove this result after replacing $X$ by
$X \setminus T$ where $T \subset Z$ is a proper closed subset.
Let $\xi$ be the generic point of $Z$.
We can choose elements $e_1, \ldots, e_{r - 1} \in \mathcal{E}_\xi$
which form part of a basis of $\mathcal{E}_\xi$.
These give rational sections $s_1, \ldots, s_{r - 1}$
of $\mathcal{O}_P(1)|_{\pi^{-1}(Z)}$ whose common zero set
is the closure of the image a rational section of
$\mathbf{P}(\mathcal{E}|_Z) \to Z$ union a closed subset whose
support maps to a proper closed subset $T$ of $Z$.
After removing $T$ from $X$ (and correspondingly $\pi^{-1}(T)$
from $P$), we see that $s_1, \ldots, s_n$ form a sequence
of global sections
$s_i \in \Gamma(\pi^{-1}(Z), \mathcal{O}_{\pi^{-1}(Z)}(1))$
whose common zero set is the image of a section $Z \to \pi^{-1}(Z)$.
Hence we see successively that
\begin{eqnarray*}
\pi^*[Z] & = & [\pi^{-1}(Z)] \\
c_1(\mathcal{O}_P(1)) \cap \pi^*[Z] & = & [Z(s_1)] \\
c_1(\mathcal{O}_P(1))^2 \cap \pi^*[Z] & = & [Z(s_1) \cap Z(s_2)] \\
\ldots & = & \ldots \\
c_1(\mathcal{O}_P(1))^{r - 1} \cap \pi^*[Z] & = &
[Z(s_1) \cap \ldots \cap Z(s_{r - 1})]
\end{eqnarray*}
by repeated applications of Lemma \ref{lemma-geometric-cap}.
Since the pushforward by $\pi$ of the image of a
section of $\pi$ over $Z$ is clearly $[Z]$ we see the result
when $\alpha = [Z]$. We omit the verification that these
arguments imply the result for a general cycle $\alpha = \sum n_j [Z_j]$.
\end{proof}

\end{document}
